\section{Introduction}

UAV multispectral imaging is increasingly used in precision agriculture and environmental monitoring, where decisions are made from cross-band products rather than from a single rendered view~\cite{colomina2014uav,turner2014uav,zhang2019uavreview,mathews2018uav}. In this setting, reconstruction quality has two coupled requirements: (1) high per-band fidelity and (2) stable cross-band geometric correspondence. The second requirement is critical for index synthesis (e.g., NDVI, GNDVI, NDRE) because even small support drift between bands can bias ratio-based measurements~\cite{rouse1974ndvi,gitelson1996gndvi,huete1988savi,goetz1985imaging,manolakis2016hyperspectral}.

Recent radiance-field and Gaussian methods greatly improved training speed and view quality~\cite{mildenhall2020nerf,barron2022mipnerf360,mueller2022instantngp,kerbl2023gaussiansplatting,yu2024mipsplatting,huang20242dgs}. Emerging multispectral and thermal Gaussian lines~\cite{meyerls2025mssplatting,thirgood2025hypergs,lu2024thermalgaussian,chen2024thermal3dgs} further show that unified spectral rendering and index synthesis are feasible. Our focus is a narrower deployment problem: UAV multispectral reconstruction under raw capture, cross-band rectification, and model-level product synthesis, where shared support must be an explicit representation contract rather than an implicit outcome.

This paper addresses that failure mode directly. We introduce \textbf{UMGS}, a UAV multispectral Gaussian splatting framework for spectral-index synthesis. UMGS formulates multispectral reconstruction as \emph{RGB-anchored geometry locking}: geometry, opacity, and Gaussian identity are learned once from RGB and then frozen during G/R/RE/NIR transfer. We evaluate two implementations of the same constraint:
\begin{itemize}
    \item \textbf{UMGS-I}: independent per-band transfer from a shared RGB anchor.
    \item \textbf{UMGS-J}: joint multispectral training with band-specific appearance banks in one unified checkpoint.
\end{itemize}
Both preserve identical support by construction; they differ only in how spectral appearance parameters are organized and optimized.

For raw UAV captures, alignment quality is another bottleneck. Cross-band FOV and resolution mismatch can make naive baselines unreliable, especially for NIR. We therefore use a baseline-aware rectification QA protocol as supporting infrastructure: it separates strict pass, pass-with-warning, and fail, and logs per-frame matching/stability diagnostics rather than relying on a single opaque gate.

Our main evaluation covers a 10-scene UAV-focused benchmark: 7 raw UAV captures and 3 aligned aerial multispectral scenes. The raw track validates end-to-end applicability (alignment, QA, reconstruction, products, and geometry diagnostics). The aligned track removes rectification confounds and supports fair common-mask comparison with MS-Splatting. Across this suite, UMGS-I and UMGS-J are near-tied and consistently preserve shared support. Scratch-MS and Unfrozen-support are useful upper-bound ablations for single-band fitting, but they do not satisfy the shared-support contract required for model-level spectral products.

The contributions are:
\begin{itemize}
    \item A geometry-locked multispectral Gaussian formulation for shared-support spectral-index synthesis.
    \item Two concrete UMGS realizations implementing the same lock with independent (UMGS-I) and joint (UMGS-J) spectral parameterizations.
    \item A UAV-focused benchmark protocol with image, index, depth-reference, cost, and external common-mask comparisons under unified reporting.
    \item Supporting raw-scene infrastructure for baseline-aware rectification QA with explicit diagnostic states and decision inputs.
\end{itemize}
